\documentclass[11pt]{article}
\usepackage{amsmath,amssymb,amsthm,mathtools,stmaryrd}
\usepackage[margin=1in]{geometry}
\usepackage{hyperref}

\newtheorem{theorem}{Theorem}
\newtheorem{lemma}{Lemma}
\newtheorem{proposition}{Proposition}
\newtheorem{corollary}{Corollary}
\theoremstyle{definition}
\newtheorem{definition}{Definition}
\newtheorem{remark}{Remark}

\newcommand{\Set}{\mathbf{Set}}
\newcommand{\Cont}{\mathbf{Cont}}
\newcommand{\ext}[1]{\llbracket #1 \rrbracket}
\newcommand{\tri}{\mathbin{\triangleleft}}
\newcommand{\pd}{\partial}
\DeclareMathOperator{\inl}{inl}
\DeclareMathOperator{\inr}{inr}

\title{A chain rule for the derivative of containers under composition}
\author{MacBeth (Kodamai / Ghani group)}
\date{2026-06-12}

\begin{document}
\maketitle

\begin{abstract}
We prove, at the level of containers over $\Set$, the chain rule
\[
  \pd(G \tri F) \;\cong\; (\pd G \tri F) \times \pd F,
\]
where $\tri$ is container composition (substitution), $\pd$ is the
one-hole-context derivative of Abbott--Altenkirch--Ghani--McBride, and $\times$
is the binary product of containers. The proof is a single explicit
isomorphism whose content is a case split on whether a chosen position lies in
the differentiated outer block or one of the untouched inner blocks. We first
dispatch the atoms and the Leibniz rule, all instances of the same one-move
template. The result is the container shadow of the Fa\`a~di~Bruno / analytic-functor
chain rule; with composition $\tri$ as the substitution operation it does not
appear in AAGM~2003, and is original in the classical container setting (cf.\
Joram--Veltri arXiv:2512.17484 for the univalent case). Verified computationally
on $23$ finite containers (explicit shape and per-shape position bijections).
\end{abstract}

\section{Setup}

We work in $\Set$ with classical logic; the one place this matters is
Lemma~\ref{lem:pointed} and is flagged there.

\begin{definition}[Container, extension]
A \emph{container} $F = S \tri P$ is a set of \emph{shapes} $S$ together with,
for each $s : S$, a set of \emph{positions} $P\,s$. Its \emph{extension} is the
functor $\ext{S \tri P} : \Set \to \Set$,
\[
  \ext{S \tri P}\,X \;=\; \textstyle\sum_{s:S}\, (P\,s \to X).
\]
A \emph{morphism} $(S \tri P) \to (S' \tri P')$ is a pair $(u, f)$ with
$u : S \to S'$ and $f : \prod_{s:S} P'(u\,s) \to P\,s$ (positions run
contravariantly). Containers and their morphisms form the category $\Cont$;
$\ext{-} : \Cont \to [\Set,\Set]$ is fully faithful. An \emph{isomorphism} of
containers is thus a bijection $u$ of shapes together with, for each $s$, a
bijection $P'(u\,s) \cong P\,s$.
\end{definition}

\begin{definition}[Operations on $\Cont$]\label{def:ops}
For $F = S \tri P$ and $G = T \tri Q$:
\begin{itemize}
  \item \emph{Sum.} $F + G = (S + T) \tri [P,Q]$, with positions inherited:
  $[P,Q](\inl s) = P\,s$, $[P,Q](\inr t) = Q\,t$. Extension $\ext{F}+\ext{G}$.
  \item \emph{Product.} $F \times G = (S \times T) \tri \lambda(s,t).\, P\,s + Q\,t$.
  Extension $\ext{F}\times\ext{G}$.
  \item \emph{Composition} (substitution) $G \tri F$, with
  $\ext{G \tri F} = \ext{G} \circ \ext{F}$:
  \[
    \text{shapes } \;\textstyle\sum_{t:T} (Q\,t \to S),
    \qquad
    \text{positions at } (t,f) \;\;\textstyle\sum_{q : Q\,t} P(f\,q).
  \]
\end{itemize}
The constants are $0 = \varnothing \tri\, !$ (no shapes), $1 = \mathbf 1 \tri (\lambda\_.\,\varnothing)$
(one shape, no positions; $\ext{1} = $ constant $\mathbf 1$), and the identity
$\mathrm{Id} = \mathbf 1 \tri (\lambda\_.\,\mathbf 1)$ (one shape, one position;
$\ext{\mathrm{Id}} = \mathrm{Id}_\Set$). For a set $A$, $K_A = A \tri (\lambda\_.\,\varnothing)$.
\end{definition}

\begin{definition}[Derivative]\label{def:deriv}
The \emph{derivative} of $F = S \tri P$ is
\[
  \pd F \;=\; \Bigl(\textstyle\sum_{s:S} P\,s\Bigr) \;\tri\;
              \lambda(s,p_0).\,\{\,p : P\,s \mid p \neq p_0\,\}.
\]
A shape of $\pd F$ is a \emph{pointed} shape $(s, p_0)$ --- a shape with a chosen
hole $p_0$ --- and its positions are the \emph{other} positions. This is the
one-hole-context functor of Abbott--Altenkirch--Ghani--McBride (``Derivatives of
Containers'', 2003/2005); $\ext{\pd F}\,X$ is the set of one-hole $F$-structures
over $X$.
\end{definition}

Throughout, for a set $B$ and a point $b_0:B$ write $B \setminus b_0 := \{b:B \mid b \neq b_0\}$.

\section{The atoms and the Leibniz rule}

Every identity below is one move: \emph{a hole in a structured position-set is a
choice of which summand the hole lives in, plus a hole there.} We record them
because the chain rule is the same move applied to a \emph{dependent} sum of
position-sets.

\begin{proposition}\label{prop:atoms}
$\pd\,\mathrm{Id} \cong 1$, \quad $\pd K_A \cong 0$ \; (in particular $\pd 1 \cong 0$).
\end{proposition}
\begin{proof}
$\pd\,\mathrm{Id}$ has shapes $\sum_{s:\mathbf 1}\mathbf 1 \cong \mathbf 1$, and at
that shape positions $\{p:\mathbf 1 \mid p\neq p_0\} = \varnothing$; so
$\pd\,\mathrm{Id} = \mathbf 1 \tri (\lambda\_.\,\varnothing) = 1$. $\pd K_A$ has
shapes $\sum_{a:A}\varnothing = \varnothing$, hence equals $0$.
\end{proof}

\begin{proposition}[Sum rule]\label{prop:sum}
$\pd(F + G) \cong \pd F + \pd G$.
\end{proposition}
\begin{proof}
Shapes of $\pd(F+G)$ are $\sum_{u : S+T}[P,Q]\,u = \sum_{s:S}P\,s \;+\; \sum_{t:T}Q\,t$,
i.e.\ shapes of $\pd F$ plus shapes of $\pd G$, by the universal property of $+$.
At $\inl(s,p_0)$ the positions are $\{p : P\,s \mid p\neq p_0\}$, exactly the
$\pd F$ positions; symmetrically at $\inr(t,q_0)$. The bijections are identities.
\end{proof}

\begin{proposition}[Leibniz / product rule]\label{prop:leibniz}
$\pd(F \times G) \cong (\pd F \times G) + (F \times \pd G)$.
\end{proposition}
\begin{proof}
A shape of $\pd(F\times G)$ is $\bigl((s,t),\,h\bigr)$ with hole
$h : P\,s + Q\,t$. The coproduct splits the hole:
\[
  \textstyle\sum_{(s,t)} (P\,s + Q\,t)
  \;\cong\; \Bigl(\sum_{(s,t)} P\,s\Bigr) + \Bigl(\sum_{(s,t)} Q\,t\Bigr)
  = \mathrm{sh}(\pd F\times G) + \mathrm{sh}(F\times\pd G).
\]
If $h = \inl p_0$, the positions $\{x : P\,s + Q\,t \mid x \neq \inl p_0\}$ are
$(P\,s \setminus p_0) + Q\,t$, exactly the positions of $\pd F \times G$ at
$((s,p_0),t)$. If $h = \inr q_0$, they are $P\,s + (Q\,t\setminus q_0)$, the
positions of $F \times \pd G$ at $(s,(t,q_0))$.
\end{proof}

\section{The chain rule}

The one ingredient beyond the atoms is the decomposition of a function out of a
pointed domain.

\begin{lemma}[Pointed-domain splitting]\label{lem:pointed}
Let $B$ be a set and $b_0 : B$. The map
$\mathbf 1 + (B\setminus b_0) \to B$, $\inl\star \mapsto b_0$,
$\inr b \mapsto b$, is a bijection. Consequently, for any set $S$,
\[
  (B \to S) \;\xrightarrow{\;\cong\;}\; S \times \bigl((B\setminus b_0) \to S\bigr),
  \qquad f \mapsto \bigl(f\,b_0,\; f|_{B\setminus b_0}\bigr).
\]
\end{lemma}
\begin{proof}
Injectivity of the first map: the images of $\inl$ and $\inr$ are disjoint
because $b_0 \notin B\setminus b_0$, and each restriction is injective.
Surjectivity: every $b:B$ satisfies $b = b_0$ or $b \neq b_0$ (excluded middle
on the proposition $b = b_0$), landing in the image of $\inl$ resp.\ $\inr$.
Applying $(-\to S)$ to $B \cong \mathbf 1 + (B\setminus b_0)$ and using
$(\mathbf 1 + C \to S) \cong S \times (C\to S)$ gives the second statement.
\end{proof}

\begin{remark}
Surjectivity is the only use of classical logic in this note. Over a general
base it requires $b_0$ to be a \emph{decidable} (complemented) point of $B$;
over $\Set$ this is automatic. This is precisely the hinge that Joram--Veltri
(arXiv:2512.17484) must handle with care in univalent foundations, where
$B\setminus b_0$ is the type $\sum_{b:B}(b \neq b_0)$ and the splitting holds
after set-truncation. For directed containers over $\Set$ it is free.
\end{remark}

\begin{theorem}[Container chain rule]\label{thm:chain}
For containers $F = S \tri P$ and $G = T \tri Q$,
\[
  \pd(G \tri F) \;\cong\; (\pd G \tri F) \times \pd F
\]
in $\Cont$.
\end{theorem}

\begin{proof}
We exhibit a shape bijection and, over it, a position bijection.

\medskip\noindent\textbf{Both sides explicitly.}
By Definitions~\ref{def:ops}--\ref{def:deriv}:

\emph{Left.} $G\tri F$ has shapes $\sum_{t:T}(Q\,t\to S)$ and positions
$R(t,f) = \sum_{q:Q t}P(f\,q)$. Hence a shape of $\pd(G\tri F)$ is
\[
  (t,\,f,\,q_0,\,p_0),\qquad
  t:T,\; f:Q\,t\to S,\; q_0:Q\,t,\; p_0:P(f\,q_0),
\]
(the hole is the position $(q_0,p_0)\in R(t,f)$), with positions
\begin{equation}\label{eq:Lpos}
  \mathrm{pos}_L = \bigl\{(q,p) : \textstyle\sum_{q:Qt}P(f\,q) \;\bigm|\; \lnot(q=q_0 \wedge p=p_0)\bigr\}.
\end{equation}

\emph{Right.} $\pd G$ has shapes $(t,q_0)$ ($t:T,\,q_0:Q\,t$) and positions
$Q\,t\setminus q_0$. So $\pd G \tri F$ has shapes
$\bigl(t,q_0,\,g\bigr)$ with $g : (Q\,t\setminus q_0)\to S$, and positions
$\sum_{q':Qt\setminus q_0} P(g\,q')$. And $\pd F$ has shapes $(s_\ast,p_\ast)$
with positions $P\,s_\ast\setminus p_\ast$. Thus a shape of the right-hand
product is
\[
  \bigl((t,q_0,g),\,(s_\ast,p_\ast)\bigr),
\]
with positions, by the product rule of Definition~\ref{def:ops},
\begin{equation}\label{eq:Rpos}
  \mathrm{pos}_R = \Bigl(\textstyle\sum_{q':Qt\setminus q_0} P(g\,q')\Bigr)
                   \;+\; \bigl(P\,s_\ast \setminus p_\ast\bigr).
\end{equation}

\medskip\noindent\textbf{Shape bijection $\varphi$.}
Send
\[
  \varphi\,(t,f,q_0,p_0) \;=\; \bigl((t,\,q_0,\,f|_{Q t\setminus q_0}),\,(f\,q_0,\,p_0)\bigr).
\]
Holding $(t,q_0)$ fixed, this is exactly the map
$f \mapsto (f\,q_0,\,f|_{Qt\setminus q_0})$ paired with the identity on $p_0$,
which is a bijection by Lemma~\ref{lem:pointed} (here $B = Q\,t$, $b_0 = q_0$):
giving $f:Q\,t\to S$ is the same as giving $s_\ast := f\,q_0 \in S$ and
$g := f|_{Qt\setminus q_0}$, and then $p_0 : P(f\,q_0) = P\,s_\ast$ is the data
$p_\ast$. Summing the bijection over all $(t,q_0)$ gives a bijection of shape
sets.

\medskip\noindent\textbf{Position bijection.}
Fix a left shape $(t,f,q_0,p_0)$ and its image, so $s_\ast = f\,q_0$,
$g = f|_{Qt\setminus q_0}$, $p_\ast = p_0$. Define
$\psi : \mathrm{pos}_L \to \mathrm{pos}_R$ on \eqref{eq:Lpos} by the hole-case split:
\[
  \psi(q,p) =
  \begin{cases}
    \inl(q,p) & q \neq q_0,\\[2pt]
    \inr(p)   & q = q_0.
  \end{cases}
\]
\emph{Well-defined.} If $q\neq q_0$ then $p : P(f\,q) = P(g\,q)$, so $(q,p)$ is a
valid element of the first summand of \eqref{eq:Rpos}. If $q = q_0$ then
$p : P(f\,q_0) = P\,s_\ast$, and the constraint $\lnot(q=q_0\wedge p=p_0)$ forces
$p \neq p_0 = p_\ast$, so $p \in P\,s_\ast\setminus p_\ast$, the second summand.

\emph{Bijective.} The two cases hit $\inl$ and $\inr$ disjointly, so it suffices
to check each is a bijection onto its summand. For the first: $q\mapsto q$,
$p\mapsto p$ with $P(f\,q)=P(g\,q)$ is the identity on
$\sum_{q\neq q_0}P(f\,q)=\sum_{q'\in Qt\setminus q_0}P(g\,q')$. For the second:
$p\mapsto p$ is the identity $P(f\,q_0)\setminus p_0 = P\,s_\ast\setminus p_\ast$.
The inverse sends $\inl(q',p')\mapsto(q',p')$ (legal since $q'\neq q_0$) and
$\inr(p')\mapsto(q_0,p')$ (legal since $p'\neq p_\ast=p_0$).

\medskip
$(\varphi,\psi)$ is a bijection of shapes together with a bijection of positions
over each shape, i.e.\ an isomorphism in $\Cont$.
\end{proof}

\paragraph{Where the two factors come from.}
The composite's position-set is the \emph{dependent} (indexed) coproduct
$\coprod_{q:Q t} P(f\,q)$ over the $G$-positions $Q\,t$. Differentiating it ---
choosing a hole --- means choosing an index $q_0$ \emph{and} a hole inside that
block. The leftover splits as
\[
  \underbrace{\coprod_{q\neq q_0} P(f\,q)}_{\text{other blocks: } \pd G\,\tri\,F}
  \;+\;
  \underbrace{\bigl(P(f\,q_0)\setminus p_0\bigr)}_{\text{this block: } \pd F}.
\]
Choosing $q_0$ and keeping the other blocks full is exactly $\pd G$ applied to
the $Q\,t$ index with the blocks substituted by $F$ --- that is $\pd G\tri F$ ---
while the punctured block at $q_0$ is exactly $\pd F$. Their juxtaposition is the
product (positions add). The chain rule is Leibniz for an \emph{indexed} product.

\section{Sanity checks}

\begin{corollary}[The $y^2$ case fixes the convention]
Take $G = y^2$, i.e.\ $T = \mathbf 1$, $Q\star = \mathbf 2$. Then
$G\tri F \cong F\times F$, and the theorem reads
$\pd(F\times F) \cong (\pd(y^2)\tri F)\times \pd F$. Since
$\pd(y^2) = \mathbf 2 \tri \mathbf 1 = y+y$, we get
$\pd(y^2)\tri F \cong F + F$, so the right side is
$(F+F)\times \pd F \cong 2\,(F\times \pd F)$. The left side is, by Leibniz,
$\pd F\times F + F\times \pd F \cong 2\,(F\times\pd F)$. They agree.
\end{corollary}

\begin{remark}[Fa\`a~di~Bruno]
Iterating Theorem~\ref{thm:chain} together with Propositions~\ref{prop:sum} and
\ref{prop:leibniz} reproduces the Fa\`a~di~Bruno formula for higher derivatives
of a composite at container level; the partitions of Fa\`a~di~Bruno are the ways
of distributing holes among the inner blocks. We do not pursue this here.
\end{remark}

\begin{remark}[Computational verification]
The explicit pair $(\varphi,\psi)$ was checked by brute force on $23$ finite
containers (one handcrafted family plus $20$ randomly generated $G,F$ with shape
counts $\le 3$ and position counts $\le 3$, including empty position-sets):
$\varphi$ is a bijection of the (finitely many) shapes, and for every left shape
$\psi$ is a bijection onto the positions of $\varphi$ of that shape. The
polynomial invariant $\#\ext{-}(n)$ agreed for $n \le 5$.
\end{remark}

\section{Status and scope}

\begin{itemize}
  \item \textbf{Theorem~\ref{thm:chain} [Original].} The chain rule for the
  derivative under \emph{composition/substitution} $\tri$ is not in
  Abbott--Altenkirch--Ghani--McBride 2003 (which give the sum and product rules
  and a derivative \emph{definition}, but not the $\tri$ chain rule). Closest
  prior art: Joram--Veltri, arXiv:2512.17484 (Dec.\ 2024, Cubical Agda), who
  prove the chain rule in univalent foundations and show the classical
  (set-truncated) case holds --- consistent with our Lemma~\ref{lem:pointed}.
  Cite them for the univalent case and for $\pd$'s definition alongside AAGM.
  \item \textbf{Derivation [MacBeth], container language.} Self-contained over
  $\Set$; the only nonconstructive step is the splitting Lemma~\ref{lem:pointed},
  flagged.
  \item \textbf{Lean [pending].} The product, coproduct and category structure
  on $\Cont$ are formalised (\texttt{Cont.lean}, PR \#12); a Lean proof of
  Theorem~\ref{thm:chain} is the natural follow-on. The proof is finite and
  bijection-based, so it should transcribe directly; the dependent transports
  around $P(f\,q_0)=P\,s_\ast$ will need the \texttt{Ext.ext\_eq} discipline noted
  for the D2/D5 laws.
  \item \textbf{Forward pointer (not chased here).} $\pd$ relates to zippers and
  to the cofree comonad; the one-hole context $\pd F$ carries a comonad structure
  whose interaction with $\tri$ deserves its own note.
\end{itemize}

\end{document}
